\section*{Chapter \ref{ch:b2:asset-prices} --- Asset Prices: A First Model}

\begin{solution}{exo:b2:asset-prices:1}
% code: exercises/solutions.py ex_state_prices
$\psi_1 + \psi_2 = 0.96$ and $2\psi_1 + 0.5\psi_2 = 1.1$. Subtracting half the first from the second,
$1.5\psi_1 = 0.62$: $\psi_1 = 0.4133$ and $\psi_2 = 0.5467$. Both are positive, so there is no
arbitrage (\cref{prop:b2:asset-prices:state-prices}).
\end{solution}

\begin{solution}{exo:b2:asset-prices:2}
% code: exercises/solutions.py ex_rf
$R_f = g^{\gamma}/\beta = 1.02^2/0.98 = 1.0616$: a real safe rate of 6.16 per cent, about five points
above what bills have paid. Patient savers who expect to be richer want to borrow against that future,
and only a high rate stops them. With a larger $\gamma$, as the \omterm{def:b2:asset-prices:equity-premium}{equity premium} needs, the gap widens.
\end{solution}

\begin{solution}{exo:b2:asset-prices:3}
% code: exercises/solutions.py ex_cumulative
Equities did best in all four countries, then bonds, then bills. French equities grew by a factor of
11.0 and American ones by 177.6. Average annual real returns: $11.0^{1/71} - 1 = 3.44$ per cent and
$177.6^{1/71} - 1 = 7.57$ per cent.
\end{solution}

\begin{solution}{exo:b2:asset-prices:4}
A rule that uses only past prices and, after trading costs, beats the return that compensates its risk,
on data the rule was not designed on, over many years. One good year proves nothing: with
unpredictable returns, about half of all rules beat the market in any given year, and one in twenty
looks significant at the 5 per cent level by chance.
\end{solution}

\begin{solution}{exo:b2:asset-prices:5}
% code: exercises/solutions.py ex_capm
$\beta_i = 0.024/0.04 = 0.6$; $\mathbb{E}[\tilde R_i] = 1 + 0.6 \times (7 - 1) = 4.6$ per cent.
\end{solution}

\begin{solution}{exo:b2:asset-prices:6}
% code: exercises/solutions.py ex_premium_size
$\mathrm{cov} = 0.5 \times 0.02 \times 0.16 = 0.0016$. Premium $0.32$ per cent for $\gamma = 2$ and
$1.6$ per cent for $\gamma = 10$. A premium of 6 per cent needs $\gamma = 0.06/0.0016 = 37.5$.
\end{solution}

\begin{solution}{exo:b2:asset-prices:7}
% code: exercises/solutions.py ex_geomean
Over the 121 years 1900--2020: United States, arithmetic mean 8.47 per cent, geometric mean 6.66; United
Kingdom, 6.63 and 4.88. The geometric mean is the constant annual return that turns a unit of 1899 into
its value in 2020, so it is the one that tells an investor what the unit became. It is smaller because
gains and losses of equal size do not cancel ($1.2 \times 0.8 < 1$); approximately, geometric mean
$\approx$ arithmetic mean $- \sigma^2/2$: $8.47 - 100 \times 0.192^2/2 = 6.63$, close to 6.66.
\end{solution}

\begin{solution}{exo:b2:asset-prices:8}
% code: exercises/solutions.py ex_minvar
Minimising the variance $\theta^2 \times 0.04 + (1 - \theta)^2 \times 0.01$ gives $\theta =
0.01/0.05 = 0.2$. The standard deviation is $\sqrt{0.2^2 \times 0.04 + 0.8^2 \times 0.01} =
\sqrt{0.008}$, that is 8.94 per cent, below either asset's.
\end{solution}

\begin{solution}{exo:b2:asset-prices:9}
Either the market misprices small firms (prices ignore information, and the \omterm{def:b2:asset-prices:excess-return}{excess return} is a
profit opportunity), or the CAPM measures their risk badly (small firms are exposed to a risk that
investors fear and \omterm{def:b2:asset-prices:beta}{beta} misses, such as doing badly in recessions). The first reading is favoured if
the \omterm{def:b2:asset-prices:excess-return}{excess return} disappears once it is published and traded on, or if it is concentrated around
announcements; the second if small firms' returns fall in bad times for the economy, beyond what their
\omterm{def:b2:asset-prices:beta}{beta} says, and the \omterm{def:b2:asset-prices:excess-return}{excess return} persists in new data and in other countries.
\end{solution}

\begin{solution}{exo:b2:asset-prices:10}
The \omterm{def:b2:insurance-and-risk-sharing:arrow-security}{Arrow security} for state 1 is a portfolio of the two securities (their payoffs are linearly
independent). Its payoff, $(1, 0)$, is never negative and positive in state 1. If its price $\psi_1$
were zero or negative, it would be an arbitrage; so no arbitrage gives $\psi_1 > 0$, and likewise
$\psi_2 > 0$. With the single security, any $\psi_2$ makes $0.5 = \psi_1 \times 1 + \psi_2 \times 0$
hold with $\psi_1 = 0.5$: every positive $\psi_2$ is consistent with no arbitrage, so \omterm{def:b2:asset-prices:state-price}{state prices} are
not determined until the market is complete.
\end{solution}

\begin{solution}{exo:b2:asset-prices:11}
Write $a = \ln\tilde g$, $b = \ln\tilde R$. From $\mathbb{E}[\tilde\zeta\tilde R] = 1$:
$\ln\beta - \gamma\,\mathbb{E}a + \mathbb{E}b + \tfrac12\mathrm{var}(b - \gamma a) = 0$. From
$\mathbb{E}[\tilde\zeta]R_f = 1$: $\ln R_f = -\ln\beta + \gamma\,\mathbb{E}a - \tfrac12\gamma^2\mathrm{var}(a)$.
Since $\ln\mathbb{E}[\tilde R] = \mathbb{E}b + \tfrac12\mathrm{var}(b)$, the first equation gives
$\ln\mathbb{E}[\tilde R] = -\ln\beta + \gamma\,\mathbb{E}a - \tfrac12\gamma^2\mathrm{var}(a) +
\gamma\,\mathrm{cov}(a, b)$. Subtracting $\ln R_f$ leaves $\gamma\,\mathrm{cov}(\ln\tilde g, \ln\tilde
R)$.
\end{solution}

\begin{solution}{exo:b2:asset-prices:12}
% code: exercises/solutions.py ex_gamma_needed, ex_gamble
(a) $\mathrm{cov}(\tilde g, \tilde R) = \rho_{gR}\,\sigma_g\sigma_R \le \sigma_g\sigma_R$, so $\gamma =
\mathbb{E}[\tilde\xi]/\mathrm{cov} \ge \mathbb{E}[\tilde\xi]/(\sigma_g\sigma_R)$. With the listing's
numbers, $0.0764/0.1629 = 0.469$ and $0.469/0.0187 = 25.1$. (b) The \omterm{def:b2:expected-utility:certainty-equivalent}{certainty equivalent} solves $c^{-24}
= 0.5 \times 0.8^{-24} + 0.5 \times 1.2^{-24}$, so $c = 0.823$ of wealth: a premium of 17.7 per cent,
almost all of the 20 per cent at stake. (c) Nobody turns down a fair coin toss for a fifth of their
wealth at that price; \omterm{def:b2:expected-utility:risk-aversion}{risk aversion} of that size contradicts the choices of \cref{ch:b2:expected-utility}
and of insurance markets, so the consumption model with CRRA utility cannot explain the premium.
\end{solution}

\begin{solution}{pb:b2:asset-prices:1}
% code: exercises/solutions.py problem
\textbf{1.} Bond: $(105, 105)$ for 100; equity: $(130, 80)$ for 95. Not proportional ($105/105 \ne
130/80$), so the two span every payoff.

\textbf{2.} $105(\psi_1 + \psi_2) = 100$ and $130\psi_1 + 80\psi_2 = 95$: $\psi_1 + \psi_2 = 0.9524$, $50\psi_1
= 95 - 80 \times 0.9524 = 18.81$, so $\psi_1 = 0.3762$ and $\psi_2 = 0.5762$.

\textbf{3.} No: both \omterm{def:b2:asset-prices:state-price}{state prices} are positive.

\textbf{4.} $R_f = 1/(\psi_1 + \psi_2) = 1/0.9524 = 1.05$.

\textbf{5.} $30 \times 0.3762 = 11.29$ dalens.

\textbf{6.} $130h_e + 105h_b = 30$ and $80h_e + 105h_b = 0$: $h_e = 0.6$, $h_b = -0.4571$ (sell
$0.4571$ bonds). Cost $0.6 \times 95 - 0.4571 \times 100 = 11.29$. Buy the option at 10 and sell the
portfolio (sell 0.6 equity, buy 0.4571 bonds): \dal{1.29} today, and the payoffs cancel in both
states.

\textbf{7.} $\zeta_1 = 0.3762/0.5 = 0.7524$, $\zeta_2 = 0.5762/0.5 = 1.1524$.

\textbf{8.} $0.5 \times 0.7524 \times 130/95 + 0.5 \times 1.1524 \times 80/95 = 0.5148 + 0.4852 = 1$.

\textbf{9.} $\mathbb{E}[\tilde R] = (0.5 \times 130 + 0.5 \times 80)/95 = 1.1053$: 10.53 per cent, a
premium of 5.53 points over the bond's 5.

\textbf{10.} $\mathbb{E}[\tilde\zeta] = 0.9524$; the deviations are $\mp 0.2$ for $\tilde\zeta$ and
$\pm 0.2632$ for $\tilde R$, so $\mathrm{cov} = -0.2 \times 0.2632 = -0.0526$, and $-1.05 \times
(-0.0526) = 0.0553$, the premium.

\textbf{11.} $\zeta_2/\zeta_1 = (0.98/1.04)^{-\gamma}$: $\gamma = \ln(1.1524/0.7524)/\ln(1.04/0.98) =
0.4263/0.0594 = 7.17$; $\beta = 0.7524 \times 1.04^{7.17} = 0.997$.

\textbf{12.} It is below the bound of 10, but only because the fund's members are assumed to swing by
6 points of consumption between the states. With the standard deviation of American consumption growth,
under 2 per cent, the same premium would need several times that $\gamma$
(\cref{exo:b2:asset-prices:12}).

\textbf{13.} Mean $0.6 \times 7 + 0.4 \times 3 = 5.4$ per cent; variance $0.36 \times 0.0324 + 0.16
\times 0.0036 + 2 \times 0.24 \times 0.00216 = 0.01328$, standard deviation 11.52 per cent.

\textbf{14.} $\theta_{\min} = (\sigma_b^2 - \sigma_{eb})/(\sigma_e^2 + \sigma_b^2 - 2\sigma_{eb}) =
(0.0036 - 0.00216)/0.03168 = 0.0455$; standard deviation 5.95 per cent.

\textbf{15.} The first-order condition $\mathbb{E}[\tilde R_e] - \mathbb{E}[\tilde R_b] =
\gamma(\theta\,0.03168 - 0.00144)$ gives $\theta^{*} = (0.04 + 0.00144\gamma)/(0.03168\gamma)$: 67.7 per
cent for $\gamma = 2$, 46.6 for $\gamma = 3$, 29.8 for $\gamma = 5$.

\textbf{16.} $0.6 \times 0.03168\gamma = 0.04 + 0.00144\gamma$: $\gamma = 0.04/0.017568 = 2.28$.

\textbf{17.} The weights are proportional to $\Sigma^{-1}(\mathbb{E}[\tilde R] - R_f)$: 42.3 per cent
equity and 57.7 bonds. Mean 4.69 per cent, standard deviation 8.97, Sharpe ratio $(4.69 - 2)/8.97 =
0.300$.

\textbf{18.} $(5.4 - 2)/11.52 = 0.295$, just below the tangency portfolio's $0.300$.

\textbf{19.} Maximising over the share $\lambda$ in the tangency portfolio: $\lambda =
(\mathbb{E}[\tilde R_t] - R_f)/(\gamma\sigma_t^2) = 0.02692/(3 \times 0.008052) = 1.11$: the fund
borrows 11 per cent of its assets at the safe rate. Its equity share overall is $1.1145 \times 0.4231 =
47.2$ per cent.

\textbf{20.} $17.55/\sqrt{60} = 2.27$ points: the mean is known only to within about 4.5 points either
way.

\textbf{21.} $t = -0.153\sqrt{60} = -1.19$, inside $\pm 1.96$: zero is not rejected.
% figdata: brenland_predict

\textbf{22.} Difference $6.28 - 5.12 = 1.16$ points; standard error $17.55\sqrt{1/38 + 1/21} = 4.77$;
$t = 0.24$. No evidence that a rise predicts anything.

\textbf{23.} About one of the twenty ($20 \times 0.05$), although no rule has any power: trying many
rules and reporting the best finds patterns that are not there.

\textbf{24.} If the return that compensates risk is higher in some periods (after a crash, say, when
investors are poorer and more averse to risk), the expected return changes predictably, and so do
realised returns on average, without any unexploited profit: the predictable part is the price of
risk. Only predictability beyond the risk-adjusted return refutes efficiency.

\textbf{25.} Without a safe asset, a fund with $\gamma = 3$ holds 46.6 per cent equity; a 60/40
portfolio is optimal for $\gamma = 2.28$; and sixty years of returns show no predictability ($t =
-1.19$ for the autocorrelation, $0.24$ for the rule).
\end{solution}
